%%% The main file. It contains definitions of basic parameters and includes all other parts.

%% Settings for single-side (simplex) printing
% Margins: left 40mm, right 25mm, top and bottom 25mm
% (but beware, LaTeX adds 1in implicitly)
\documentclass[12pt,a4paper]{report}
\setlength\textwidth{145mm}
\setlength\textheight{247mm}
\setlength\oddsidemargin{15mm}
\setlength\evensidemargin{15mm}
\setlength\topmargin{0mm}
\setlength\headsep{0mm}
\setlength\headheight{0mm}
% \openright makes the following text appear on a right-hand page
\let\openright=\clearpage

%% Settings for two-sided (duplex) printing
% \documentclass[12pt,a4paper,twoside,openright]{report}
% \setlength\textwidth{145mm}
% \setlength\textheight{247mm}
% \setlength\oddsidemargin{14.2mm}
% \setlength\evensidemargin{0mm}
% \setlength\topmargin{0mm}
% \setlength\headsep{0mm}
% \setlength\headheight{0mm}
% \let\openright=\cleardoublepage

%% Generate PDF/A-2u
\usepackage[a-2u]{pdfx}

%% Character encoding: usually latin2, cp1250 or utf8:
\usepackage[utf8]{inputenc}

%% Prefer Latin Modern fonts
\usepackage{lmodern}

%% Further useful packages (included in most LaTeX distributions)
\usepackage{amsmath}        % extensions for typesetting of math
\usepackage{amsfonts}       % math fonts
\usepackage{amsthm}         % theorems, definitions, etc.
\usepackage{bbding}         % various symbols (squares, asterisks, scissors, ...)
\usepackage{bm}             % boldface symbols (\bm)
\usepackage{graphicx}       % embedding of pictures
\usepackage{fancyvrb}       % improved verbatim environment
\usepackage{natbib}         % citation style AUTHOR (YEAR), or AUTHOR [NUMBER]
\usepackage[nottoc]{tocbibind} % makes sure that bibliography and the lists
			    % of figures/tables are included in the table
			    % of contents
\usepackage{dcolumn}        % improved alignment of table columns
\usepackage{booktabs}       % improved horizontal lines in tables
\usepackage{paralist}       % improved enumerate and itemize
\usepackage{xcolor}         % typesetting in color

%%% Basic information on the thesis

% Thesis title in English (exactly as in the formal assignment)
\def\ThesisTitle{Magnetic Transport Along Translationally Invariant Obstacles}

% Author of the thesis
\def\ThesisAuthor{Michal Grňo}

% Year when the thesis is submitted
\def\YearSubmitted{2021}

% Name of the department or institute, where the work was officially assigned
% (according to the Organizational Structure of MFF UK in English,
% or a full name of a department outside MFF)
\def\Department{Institute of Theoretical Physics}

% Is it a department (katedra), or an institute (ústav)?
\def\DeptType{institute}

% Thesis supervisor: name, surname and titles
\def\Supervisor{prof. RNDr. Pavel Exner, DrSc.}

% Supervisor's department (again according to Organizational structure of MFF)
\def\SupervisorsDepartment{Institute of Theoretical Physics}

% Study programme and specialization
\def\StudyProgramme{Fyzika}
\def\StudyBranch{FOF}

% An optional dedication: you can thank whomever you wish (your supervisor,
% consultant, a person who lent the software, etc.)
\def\Dedication{%
Dedication.
}

% Abstract (recommended length around 80-200 words; this is not a copy of your thesis assignment!)
\def\Abstract{%
Abstract.
}

% 3 to 5 keywords (recommended), each enclosed in curly braces
\def\Keywords{%
{two-dimensional quantum systems}, {magnetic field}, {nonlocal perturbations}, {translational invariance}, {spectral properties}
}

%% The hyperref package for clickable links in PDF and also for storing
%% metadata to PDF (including the table of contents).
%% Most settings are pre-set by the pdfx package.
\hypersetup{unicode}
\hypersetup{breaklinks=true}

% Definitions of macros (see description inside)
%%% This file contains definitions of various useful macros and environments %%%
%%% Please add more macros here instead of cluttering other files with them. %%%

%%% Minor tweaks of style

% These macros employ a little dirty trick to convince LaTeX to typeset
% chapter headings sanely, without lots of empty space above them.
% Feel free to ignore.
\makeatletter
\def\@makechapterhead#1{
  {\parindent \z@ \raggedright \normalfont
   \Huge\bfseries \thechapter. #1
   \par\nobreak
   \vskip 20\p@
}}
\def\@makeschapterhead#1{
  {\parindent \z@ \raggedright \normalfont
   \Huge\bfseries #1
   \par\nobreak
   \vskip 20\p@
}}
\makeatother

% This macro defines a chapter, which is not numbered, but is included
% in the table of contents.
\def\chapwithtoc#1{
\chapter*{#1}
\addcontentsline{toc}{chapter}{#1}
}

% Draw black "slugs" whenever a line overflows, so that we can spot it easily.
\overfullrule=1mm

%%% Macros for definitions, theorems, claims, examples, ... (requires amsthm package)

\theoremstyle{plain}
\newtheorem{thm}{Theorem}
\newtheorem{lemma}[thm]{Lemma}
\newtheorem{claim}[thm]{Claim}
\newtheorem{defn}[thm]{Definition}

\theoremstyle{remark}
\newtheorem*{cor}{Corollary}
\newtheorem*{rem}{Remark}
\newtheorem*{example}{Example}

%%% An environment for proofs

\newenvironment{myproof}{
  \par\medskip\noindent
  \textit{Proof}.
}{
\newline
\rightline{$\qedsymbol$}
}

%%% An environment for typesetting of program code and input/output
%%% of programs. (Requires the fancyvrb package -- fancy verbatim.)

\DefineVerbatimEnvironment{code}{Verbatim}{fontsize=\small, frame=single}

%%% The field of all real and natural numbers
\renewcommand{\C}{\mathbb{C}}
\newcommand{\K}{\mathbb{K}}
\newcommand{\R}{\mathbb{R}}
\newcommand{\Z}{\mathbb{Z}}
\newcommand{\N}{\mathbb{N}}

%%% Useful operators for statistics and probability
\DeclareMathOperator{\pr}{\textsf{P}}
\DeclareMathOperator{\E}{\textsf{E}\,}
\DeclareMathOperator{\var}{\textrm{var}}
\DeclareMathOperator{\sd}{\textrm{sd}}

%%% Transposition of a vector/matrix
\newcommand{\T}[1]{#1^\top}

%%% Various math goodies
\newcommand{\goto}{\rightarrow}
\newcommand{\gotop}{\stackrel{P}{\longrightarrow}}
\newcommand{\maon}[1]{o(n^{#1})}
\newcommand{\abs}[1]{\left|{#1}\right|}
\newcommand{\isqr}[1]{\frac{1}{\sqrt{#1}}}

%%% Various table goodies
\newcommand{\pulrad}[1]{\raisebox{1.5ex}[0pt]{#1}}
\newcommand{\mc}[1]{\multicolumn{1}{c}{#1}}


%%% Custom commands
\renewcommand{\d}[1]{\, \mathrm{d}#1 \,}
\renewcommand{\i}{\mathrm{i}}
\newcommand{\Cont}{C}
\newcommand{\Hilb}{\mathcal{H}}
\newcommand{\Schwa}{\mathcal{S}}
\newcommand{\Domain}{\mathrm{D}}
\newcommand{\Fourier}{\mathcal{F}}
\newcommand{\Nbhood}{\mathcal{U}}
\newcommand{\norm}[1]{\big\lVert #1 \big\rVert}
\newcommand{\Hf}{\mathscr{H}}
\newcommand{\Af}{\mathscr{A}}
\newcommand{\ve}{\varepsilon}
\newcommand{\hypF}{\vphantom{\mathrm{F}}_1\hspace{-0.7pt}\mathrm{F}_{\!1}}
\newcommand{\supp}{\operatorname{supp}}

\newcommand{\Sp}{\sigma}
\newcommand{\SpP}{\sigma_{\mathrm{p}}}
\newcommand{\SpAc}{\sigma_{\mathrm{ac}}}
\newcommand{\SpSc}{\sigma_{\mathrm{sc}}}
\newcommand{\SpD}{\sigma_{\mathrm{disc}}}
\newcommand{\SpE}{\sigma_{\mathrm{ess}}}

\newcommand{\tildeA}{\overset{\;\,\scriptscriptstyle\sim}{A}}
\newcommand{\tildeH}{\overset{\;\,\scriptscriptstyle\sim}{H}}


\def\ph{\phantom}
\def\vph{\vphantom}
\def\hph{\hphantom}
\def\rzw{\mathrlap}
\def\lzw{\mathllap}
\def\czw{\mathclap}
\newcommand*{\mask}[2]{\mathord{\makebox[\widthof{\(#1\)}]{\(#2\)}}}
\newcommand{\scalemath}[2]{\scalebox{#1}{\mbox{\ensuremath{\displaystyle #2}}}}


\newcommand{\dd}[2]{\frac{\mathrm{d} #1}{\mathrm{d} #2}}
\newcommand{\pd}[2]{\frac{\partial  #1}{\partial  #2}}
\newcommand{\e}[1]{\mathrm{e}^{#1}}

\newcommand{\coloneqq}{:=}
\newcommand{\eqqcolon}{=:}

\newcommand{\digamma}{\Psi} % haha dobrý

\newcommand\numberthis{\addtocounter{equation}{1}\tag{\theequation}}



% Title page and various mandatory informational pages
\begin{document}
%%% Title page of the thesis and other mandatory pages

%%% Title page of the thesis

\pagestyle{empty}
\hypersetup{pageanchor=false}
\begin{center}

\centerline{\mbox{\includegraphics[width=166mm]{../img/logo-en.pdf}}}

\vspace{-8mm}
\vfill

{\bf\Large BACHELOR THESIS}

\vfill

{\LARGE\ThesisAuthor}

\vspace{15mm}

{\LARGE\bfseries\ThesisTitle}

\vfill

\Department

\vfill

{
\centerline{\vbox{\halign{\hbox to 0.45\hsize{\hfil #}&\hskip 0.5em\parbox[t]{0.45\hsize}{\raggedright #}\cr
Supervisor of the bachelor thesis:&\Supervisor \cr
\noalign{\vspace{2mm}}
Study programme:&\StudyProgramme \cr
\noalign{\vspace{2mm}}
Study branch:&\StudyBranch \cr
}}}}

\vfill

% Zde doplňte rok
Prague \YearSubmitted

\end{center}

\newpage

%%% Here should be a bound sheet included -- a signed copy of the "bachelor
%%% thesis assignment". This assignment is NOT a part of the electronic
%%% version of the thesis. DO NOT SCAN.

%%% A page with a solemn declaration to the bachelor thesis

\openright
\hypersetup{pageanchor=true}
\pagestyle{plain}
\pagenumbering{roman}
\vglue 0pt plus 1fill

\noindent
I declare that I carried out this bachelor thesis independently, and only with the cited
sources, literature and other professional sources. It has not been used to obtain another
or the same degree.

\medskip\noindent
I understand that my work relates to the rights and obligations under the Act No.~121/2000 Sb.,
the Copyright Act, as amended, in particular the fact that the Charles
University has the right to conclude a license agreement on the use of this
work as a school work pursuant to Section 60 subsection 1 of the Copyright~Act.

\vspace{10mm}

\hbox{\hbox to 0.5\hsize{%
In \hbox to 6em{\dotfill} date \hbox to 6em{\dotfill}
\hss}\hbox to 0.5\hsize{\dotfill\quad}}
\smallskip
\hbox{\hbox to 0.5\hsize{}\hbox to 0.5\hsize{\hfil Author's signature\hfil}}

\vspace{20mm}
\newpage

%%% Dedication

\openright

\noindent
\Dedication

\newpage

%%% Mandatory information page of the thesis

\openright

\vbox to 0.5\vsize{
\setlength\parindent{0mm}
\setlength\parskip{5mm}

Title:
\ThesisTitle

Author:
\ThesisAuthor

\DeptType:
\Department

Supervisor:
\Supervisor, \SupervisorsDepartment

Abstract:
\Abstract

Keywords:
\Keywords

\vss}

\newpage

\openright
\pagestyle{plain}
\pagenumbering{arabic}
\setcounter{page}{1}


%%% A page with automatically generated table of contents of the bachelor thesis

\tableofcontents

%%% Each chapter is kept in a separate file
\chapter*{List of symbols}
\addcontentsline{toc}{chapter}{List of symbols}

{\renewcommand{\arraystretch}{1.3}
\begin{table}[h!]
    \begin{tabular}{c|p{0.7\linewidth}}
        \makecell[tc]{$\Cont^k(\Omega, \K)$, \\ $\Cont(\Omega, \K)$} & The space of functions $\Omega \subseteq \R \to \K$ with $k$ continuous derivatives. The space of continuous functions. \\
        $\Cont^\infty_0(\Omega, \K)$ & The space of $C^\infty$ functions with compact support in $\Omega$. \\
        $\Domain(T)$ & The domain of operator $T$, usually dense in $\Hilb$. \\
        $D_\nu(w)$ & The parabolic cylinder function. \\
        $\Fourier$ & The Fourier-Plancherel operator on $L^2(\R)$. \\
        $\hypF(\alpha, \beta; z)$ & The confluent hypergeometric function of the first kind. \\
        $\Hilb$ & A separable Hilbert space. \\
        $H, \; \Hf(\xi)$ & A Hamiltonian operator; a fibre of the Hamiltonian. \\
        $H_n(x)$ & The $n$-th Hermite polynomial. \\
        $L^p(M, \d{\mu}, V)$ & The space of $p$-integrable functions from measure space $(M, \mu)$ to vector space $V$. Specifically for $p=2$, a Hilbert space with inner product $(\psi, \phi)_{L^2} = \int_M (\psi, \phi)_V \d{\mu}$. \\
        $L^p(\Omega)$ & As above, but $M=\Omega \subseteq \R^N$, $\mu$ is the Lebesgue measure and $V=\C$. \\
        $L^p_\rho(\Omega)$ & A weighted Lebesgue space, a shorthand for $L^p(\Omega \subseteq \R^N, \d{\mu}, \C)$, where $\rho$ is a real function and $\mu(M) = \int_M \rho \d{\lambda}$ is a rescaling of the Lebesgue measure $\lambda$. \\
        $L^1_{\mathrm{loc}}(\Omega)$ & The space of functions that are $L^1(K)$ for every compact $K \subset \Omega$. \\
        $\N, \; \N_0$ & The set of positive integers; the set of non-negative integers. \\
        $\vv P, P_x, P_y, P_z$ & Momentum operator – a self-adjoint operator, such that $P_x \, f(x,...) = -\i \pd{}{x} f(x,...) \: .$ \\
        $\vv Q, Q_x, Q_y, Q_z$ & Position operator – a self-adjoint operator, such that $Q_x \, f(x,...) = x \, f(x,...) \: .$ \\
        $W^{k,p}(\Omega)$ & The Sobolev space – the space of integrable functions $f$, such that $f^{(\alpha)} \in L^p(\Omega)$, where $\alpha$ is a multi-index and $|\alpha| \leq k$. \\
        $\Gamma(z), \digamma(z)$ & The gamma function, the digamma function. \\
        $\mu$ & A $\sigma$-finite measure, usually the Lebesgue measure. \\
        \makecell[tc]{$\Sp(T), \SpP(T),$\\$\SpAc(T), \SpSc(T),$\\$\SpD(T), \SpE(T)$} & The spectrum of normal operator $T$; the point, absolutely continuous, singular continuous, discrete, essential spectrum of $T$. $\Sp(T) = \SpP \cup \SpAc \cup \SpSc = \SpD \cup \SpE \:$. \\
        $\nabla, \; \nabla \times, \; \Delta$ & Gradient, rotation, Laplace operator. \\
        $\Delta^\Omega_{\mathrm{D}}, \Delta^\Omega_{\mathrm{D},A}$ & The Dirichlet Laplacian, defined on functions from $L^2(\Omega)$ with a Dirichlet boundary condition; a “magnetic” Dirichlet Laplacian given by the vector potential $A$.
    \end{tabular}
\end{table}
}

\chapter*{Introduction}
\addcontentsline{toc}{chapter}{Introduction}


\chapter{Title of the first chapter}

An~example citation: \cite{Andel07}

\section{Title of the first subchapter of the first chapter}

\section{Title of the second subchapter of the first chapter}

\chapter{Title of the second chapter}

\section{Title of the first subchapter of the second chapter}

\section{Title of the second subchapter of the second chapter}


\chapter*{Conclusion}
\addcontentsline{toc}{chapter}{Conclusion}

This thesis aimed to summarize the current knowledge of translationally invariant magnetic Hamiltonians and complement that knowledge with two previously unstudied systems.

One of the systems that we studied was the Landau Hamiltonian with a Dirac delta-interaction of strength $\alpha$ supported on a line. We verified that the Hamiltonian is self-adjoint and bounded from below, that it is a fibred operator and the spectrum of its fibre is discrete and nondegenerate. Then we found implicit expressions for the eigenvalues and eigenfunctions of the fibre Hamiltonian. Finally, we proved that: for a repulsive interaction ($\alpha>0$), there are spectral bands that start on each Landau level and extend up, for an attractive interaction ($\alpha < 0$), the spectral bands start on each Landau level and extend down. In both cases, the spectral bands never merge – there is a forbidden energy gap between each two neighbouring Landau levels. The spectrum of the Hamiltonian is absolute continuous; its singular continuous and pure point spectra are empty.

The other system that we examined was the Landau Hamiltonian on a half-plane with a Robin boundary condition. As with the delta-interaction, we verified the self-adjointness and boundedness from below, and that the fibre Hamiltonian has a discrete nondegenerate spectrum. Then we used the properties of the fibre to prove that: the spectrum of the Hamiltonian is absolute continuous, starting from a minimum below the first Landau level and extending up ad infinitum. We have numerically calculated estimates of the minimum energy for different values of the boundary parameter $\alpha$.

By comparing the (numerically computed) eigenvalues of the two systems' fibre Hamiltonians, we observed that if $p_y \ll 0$ and $|\alpha|$ is large enough, then the energy levels of the Robin-boundary fibre Hamiltonian with parameter $\alpha$ are a good approximation of the fibre Hamiltonian of a $\delta$-interaction of strength $-\alpha$.

Topics for further studies on the spectral properties of translationally invariant magnetic Hamiltonians are far from exceeded. Regarding the two Hamiltonians we studied, future work might investigate the stability of their spectra under perturbations by various classes of disorder potentials, or find upper and lower estimates for the minimum energy of the Robin half-plane Hamiltonian. Other Landau Hamiltonians with potential perturbations which, to the author's knowledge, are yet to be studied include: multiple $\delta$-interactions on parallel lines, a $\delta'$-interaction\footnote{The plausibility of a $\delta'$-interaction in one dimension is discussed in Chapter III.3 of \cite{Albeverio2005}.} supported on a line or multiple parallel lines, and a strip with Neumann or Robin boundaries.


%%% Bibliography
%%% Bibliography (literature used as a source)
%%%
%%% We employ bibTeX to construct the bibliography. It processes
%%% citations in the text (e.g., the \cite{...} macro) and looks up
%%% relevant entries in the bibliography.bib file.
%%%
%%% The \bibliographystyle command selects, which style will be used
%%% for references from the text. The argument in curly brackets is
%%% the name of the corresponding style file (*.bst). Both styles
%%% mentioned in this template are included in LaTeX distributions.

\bibliographystyle{plainnat}    %% Author (year)
% \bibliographystyle{unsrt}     %% [number]

\renewcommand{\bibname}{Bibliography}

%%% Generate the bibliography. Beware that if you cited no works,
%%% the empty list will be omitted completely.

\bibliography{bibliography}

%%% If case you prefer to write the bibliography manually (without bibTeX),
%%% you can use the following. Please follow the ISO 690 standard and
%%% citation conventions of your field of research.

% \begin{thebibliography}{99}
%
% \bibitem{lamport94}
%   {\sc Lamport,} Leslie.
%   \emph{\LaTeX: A Document Preparation System}.
%   2nd edition.
%   Massachusetts: Addison Wesley, 1994.
%   ISBN 0-201-52983-1.
%
% \end{thebibliography}


%%% Figures used in the thesis (consider if this is needed)
\listoffigures

%%% Tables used in the thesis (consider if this is needed)
%%% In mathematical theses, it could be better to move the list of tables to the beginning of the thesis.
\listoftables

%%% Abbreviations used in the thesis, if any, including their explanation
%%% In mathematical theses, it could be better to move the list of abbreviations to the beginning of the thesis.
\chapwithtoc{List of Abbreviations}

%%% Attachments to the bachelor thesis, if any. Each attachment must be
%%% referred to at least once from the text of the thesis. Attachments
%%% are numbered.
%%%
%%% The printed version should preferably contain attachments, which can be
%%% read (additional tables and charts, supplementary text, examples of
%%% program output, etc.). The electronic version is more suited for attachments
%%% which will likely be used in an electronic form rather than read (program
%%% source code, data files, interactive charts, etc.). Electronic attachments
%%% should be uploaded to SIS and optionally also included in the thesis on a~CD/DVD.
%%% Allowed file formats are specified in provision of the rector no. 72/2017.
\appendix
\chapter{Attachments}

\section{First Attachment}

\openright
\end{document}
